\section{Scaling rigidity of the Cayley collision family}
\label{sec:cayley-scaling}
Consider the off-diagonal collision polynomial
\[
 G_{\rho,\sigma}(U,V)=\sigma(\rho-1)U^2V^2+\rho(1-\sigma)U^2
          +2(\sigma-\rho)UV+(1-\sigma)V^2+\rho-1.
\]
Assume characteristic different from two, $\rho,\sigma\notin\{0,1,-1\}$,
and $\rho\ne\sigma$. These assumptions hold for distinct nonunit parameters
in an odd-order multiplicative subgroup.

\begin{proposition}
Among this family a scalar and diagonal-coordinate scaling class has exactly
the two members $(\rho,\sigma)$ and $(\rho^{-1},\sigma^{-1})$ over an
algebraic closure. Allowing invertible Laurent monomial coordinate changes
and Laurent monomial multipliers enlarges it to at most four members, adding
$(\sigma,\rho)$ and $(\sigma^{-1},\rho^{-1})$.
For parameters in an odd-order subgroup all these equivalences can be
realized by transformations preserving its Cartesian square.
\end{proposition}
\begin{proof}
For coefficients $(a,b,c,d,e)$ of $(U^2V^2,U^2,UV,V^2,1)$, the quantities
$I=ae/(bd)$ and $J=c^2/(bd)$ are scaling invariants. Here
\[
 I=\frac{\sigma(\rho-1)^2}{\rho(\sigma-1)^2},\qquad
 J=\frac{4(\sigma-\rho)^2}{\rho(\sigma-1)^2}.
\]
Writing $z=(\sigma-\rho)/(\sigma-1)$ and $c_0=1-I+J/4$, direct algebra gives
\[
 z^2=J\rho/4,\quad c_0\rho=(\rho+1)z,\quad
                       J(\rho+1)^2=4c_0^2\rho.
\]
The last equation is a reciprocal quadratic with $J\ne0$, so it allows at
most two values of $\rho$. Each determines $z$ and then
$\sigma=(z-\rho)/(z-1)$. Simultaneous inversion preserves the invariants
and gives the second distinct pair. If
$\alpha^2=(\rho\sigma)^{-1}$ and $\beta=\rho\alpha$, then
\[
 G_{\rho^{-1},\sigma^{-1}}(U,V)
     =-\rho^{-1}G_{\rho,\sigma}(\alpha U,\beta V).
\]
Squaring is bijective in an odd-order subgroup, proving that restricted
statement as well.

The Newton polygon is the square $[0,2]^2$ with its centre as the only other
support point. Its affine unimodular symmetries have signed-permutation
linear parts. Interchanging the coordinates preserves $(I,J)$; inverting one
coordinate changes it to $(I^{-1},J/I)$, the invariant pair obtained by
interchanging $\rho,\sigma$. The first part therefore bounds the class by
four, and the identity
$G_{\sigma,\rho}(U,V)=-V^2G_{\rho,\sigma}(U,V^{-1})$ realizes the swap.
There are only two distinct members when $\rho\sigma=1$, and four otherwise.
\end{proof}

This rules out the proposed direct conversion of an unbounded two-parameter
collision family into rescalings of one fixed polynomial by these coordinate
changes. It does not rule out other curve-averaging methods, and it supplies
no improvement of the active $K^{8/3}$ pointwise energy exponent.

\section{A global Cayley subgroup-curve energy bound}
\label{sec:cayley-curve-uniform}
The scaling-rigidity result above explains why the two-parameter collision
family cannot be reduced to one polynomial by elementary coordinate changes.
There is nevertheless a uniform bound for every subgroup size, obtained by
applying the positive-characteristic curve theorem of Corvaja--Zannier~\cite{CorvajaZannier2013},
\emph{J. Eur. Math. Soc.} \textbf{15} (2013), Theorem~2/Corollary~2,
to each off-diagonal ratio curve.  The exponent remains
$8/3$, but the earlier parity and size restrictions disappear.

Let $p$ be an odd prime, let $H\leq\mathbf F_p^*$ have order $K\geq2$, and put
\[
 c(h)=\frac{h+1}{h-1},\qquad A=c(H\setminus\{1\}),
\]
with ratio-fibre map
\[
 d_\rho(x)=\frac{(\rho x-1)(x-1)}{(1-\rho)x},
 \qquad \rho\in H\setminus\{1\}.
\]
For $\rho,\sigma\in H\setminus\{1\}$ define the extended collision count
\[
 N(\rho,\sigma)=\#\{(x,y)\in H^2:d_\rho(x)=d_\sigma(y)\}.
\]
The original pair-key energy satisfies
\begin{equation}\label{eq:czu-ec}
 E_C:=\sum_t C_H(t)^2=E_+(A)-(K-1)^2.
\end{equation}

\begin{theorem}[Uniform Cayley curve bound]\label{thm:czu-cayley}
For every odd prime $p$, every subgroup $H\leq\mathbf F_p^*$ of order
$K\geq2$, and every distinct $\rho,\sigma\in H\setminus\{1\}$,
\begin{equation}\label{eq:czu-point}
 N(\rho,\sigma)\leq
 \min\!\left\{2K,
 \max\!\left(3\sqrt[3]{32}\,K^{2/3},\frac{48K^2}{p}\right)\right\}.
\end{equation}
If
\[
 B(K,p)=\min\!\left\{2K,
 \max\!\left(3\sqrt[3]{32}\,K^{2/3},\frac{48K^2}{p}\right)\right\}
\]
and
\[
 D_K=2K^2-7K+4+\gcd(2,K),
\]
then
\begin{equation}\label{eq:czu-energy-exact}
 E_+(A)-(K-1)^2
 \leq D_K+(K-1)(K-2)\bigl(B(K,p)-4\bigr).
\end{equation}
Consequently, uniformly for every subgroup size,
\begin{equation}\label{eq:czu-energy-asymp}
 E_+(A)\ll K^{8/3}+\frac{K^4}{p},
\end{equation}
with an absolute implied constant.
\end{theorem}

\begin{proof}
Clearing denominators in $d_\rho(x)=d_\sigma(y)$ gives
\begin{equation}\label{eq:czu-F}
 F_{\rho,\sigma}(x,y)
 =(1-\sigma)y(\rho x-1)(x-1)
 -(1-\rho)x(\sigma y-1)(y-1).
\end{equation}
For $\rho\neq\sigma$ this has bidegree $(2,2)$.  Viewed as a quadratic in
$y$, its discriminant is a quartic in $x$, and direct elimination gives
\begin{equation}\label{eq:czu-disc}
 \operatorname{disc}_x\!\bigl(\operatorname{disc}_yF_{\rho,\sigma}\bigr)
 =4096\rho^2\sigma^2(\rho-1)^6(\rho-\sigma)^2(\sigma-1)^6.
\end{equation}
The right side is nonzero in odd characteristic for distinct nonzero
nonunit parameters.  Hence the quadratic discriminant is squarefree and is
not a square in $\overline{\mathbf F}_p(x)$; therefore the curve is
geometrically irreducible.

Its closure in $\mathbf P^1\times\mathbf P^1$ has exactly the four boundary
points
\[
 (0,0),\quad(0,\infty),\quad(\infty,0),\quad(\infty,\infty).
\]
In the corresponding affine charts the pairs of first derivatives are
\[
 (\rho-1,1-\sigma),\quad
 (\sigma(\rho-1),1-\sigma),\quad
 (\rho-1,\rho(1-\sigma)),\quad
 (\sigma(\rho-1),\rho(1-\sigma)),
\]
so none is singular.  Equivalently, \eqref{eq:czu-disc} exhibits the
normalization as a double cover of $\mathbf P^1$ branched at four points,
thus of genus one.  Since a bidegree $(2,2)$ divisor has arithmetic genus
one, there are no hidden projective singularities.

The coordinate functions $x$ and $y$ both have degree two.  Their valuation
pairs at the four boundary points are
\[
 (1,1),\quad(1,-1),\quad(-1,1),\quad(-1,-1),
\]
which proves multiplicative independence modulo constants.  With $S$ the
boundary divisor, the affine Euler characteristic entering
Corvaja--Zannier is therefore
\[
 \chi=|S|+2g-2=4.
\]
The degree-two coordinate maps are separable because $p$ is odd, while
$p\nmid K$ because $K\mid p-1$.  Corollary~2 of that paper~\cite{CorvajaZannier2013}, applied with
$m_1=m_2=K$, $d_1=d_2=2$, and $\chi=4$, now yields
\[
 N(\rho,\sigma)
 \leq\max\!\left(3\sqrt[3]{32}\,K^{2/3},\frac{48K^2}{p}\right).
\]
For each fixed $x\in H$, equation \eqref{eq:czu-F} is quadratic in $y$,
which also gives the trivial bound $N(\rho,\sigma)\leq2K$.  This proves
\eqref{eq:czu-point}.

It remains to restore the exact pair-key domain.  Extending one ratio fibre
to all $x\in H$ adds precisely $x=1$ and $x=\rho^{-1}$, and both have value
zero.  Hence every ordered parameter pair contributes exactly four
artificial zero--zero collisions, giving the exact identity
\begin{equation}\label{eq:czu-restore}
 E_C=\sum_{\rho,\sigma\neq1}N(\rho,\sigma)-4(K-1)^2.
\end{equation}
On the diagonal,
\[
 F_{\rho,\rho}(x,y)=-(\rho-1)(x-y)(\rho xy-1).
\]
Summing the intersections of these two components over $\rho\neq1$ gives
$K-\gcd(2,K)$ and therefore
\[
 \sum_{\rho\neq1}N(\rho,\rho)
 =2K(K-1)-K+\gcd(2,K).
\]
After subtracting the four artificial collisions for every diagonal ratio,
the total diagonal contribution to $E_C$ is exactly $D_K$.  There are
$(K-1)(K-2)$ ordered off-diagonal parameter pairs; inserting
\eqref{eq:czu-point} into \eqref{eq:czu-restore} proves
\eqref{eq:czu-energy-exact}.  Finally
$B(K,p)\ll K^{2/3}+K^2/p$, and the diagonal term is $O(K^2)$, which gives
\eqref{eq:czu-energy-asymp}.
\end{proof}

\begin{remark}[Range and sharp frontier]
The characteristic-free term dominates whenever
$128K^4\leq p^3$, equivalently
$K\leq2^{-7/4}p^{3/4}$.  Its coefficient is
$3\sqrt[3]{32}=9.5244063118\ldots$.  This removes the old odd-order and
subgroup-size restrictions from the project curve route, but it does not
improve the exponent.  The centered pointwise target
\[
 E_+(A)=\frac{(K-1)^4}{p}+O\!\left(K^{2+o(1)}\right)
\]
remains open and requires cancellation across the two-parameter family,
not merely a worst-curve estimate.
\end{remark}

\begin{remark}[Exact regression boundary]
The maintained exact verifier checks the discriminant identity, all four
boundary gradients, diagonal factorization, restoration correction and the
integer form of the energy estimate for every subgroup order $K\geq3$
occurring for odd primes $p\leq257$.  This is a finite algebra and
implementation regression; it does not re-prove the external Corvaja--Zannier theorem.
\end{remark}


\subsection{Dense-regime Fourier complement}
\label{subsec:cayley-fourier-hybrid}
The genus-one curve estimate is strongest in the sparse and intermediate
regimes.  A different argument gives a centered estimate for dense subgroups.
Let $m=(p-1)/K$, $n=K-1$, and use the unnormalised transform
\[
 \widehat{1_A}(q)=\sum_{a\in A}e_p(qa).
\]
If $S=H^\perp$ is the group of multiplicative characters trivial on $H$,
then $|S|=m$ and for $q\ne0$
\[
 \widehat{1_A}(q)=\frac1m\sum_{\chi\in S}T_\chi(q),\qquad
 T_\chi(q)=\sum_{h\in\mathbf F_p^*,\ h\ne1}
     \chi(h)e_p\!\left(q\frac{h+1}{h-1}\right).
\]
The trivial character contributes exactly $-e_p(q)-e_p(-q)$.  For
$\chi\ne1$, the substitution $z=(h+1)/(h-1)$ gives the mixed
multiplicative--additive sum
\[
 T_\chi(q)=\sum_{z\in\mathbf F_p\setminus\{-1,1\}}
 \chi\!\left(\frac{z+1}{z-1}\right)e_p(qz).
\]
The classical mixed Weil bound gives $|T_\chi(q)|\le2\sqrt p$; the
Kummer factor has one simple zero and one simple pole and the additive
factor has its sole pole at infinity. Consequently
\begin{equation}\label{eq:cayley-fourier-max}
 |\widehat{1_A}(q)|\le
 M_m(p):=\frac{2((m-1)\sqrt p+1)}m\le2\sqrt p .
\end{equation}

\begin{theorem}[Dense centered Cayley energy bound]
For every odd prime $p$ and every subgroup $H\le\mathbf F_p^*$ of order
$K\ge2$, with $A=c(H\setminus\{1\})$, $n=K-1$ and
$m=(p-1)/K$,
\begin{align}
0\le E_+(A)-\frac{n^4}{p}
&\le \frac{M_m(p)^2}{p}\,n(p-n) \\
&=4\frac{((m-1)\sqrt p+1)^2}{m^2p}\,n(p-n) \\
&\le4n(p-n).                                      \label{eq:cayley-fourier-energy}
\end{align}
Combining this with Theorem~\ref{thm:czu-cayley},
\[
 E_+(A)-\frac{(K-1)^4}{p}
 \ll \min\!\left\{pK,\ K^{8/3}+\frac{K^4}{p}\right\}.
\]
Thus if $K=p^{\alpha+o(1)}$ with $\alpha>3/5$, the centered error has
strictly smaller exponent than $K^{8/3}$; in $K$-notation the Fourier
branch is $K^{1+1/\alpha+o(1)}$.  For example $\alpha=2/3$ gives a
$K^{5/2+o(1)}$ centered bound.
\end{theorem}
\begin{proof}
Parseval gives
\[
 \sum_q|\widehat{1_A}(q)|^2=pn,\qquad
 \sum_{q\ne0}|\widehat{1_A}(q)|^2=n(p-n).
\]
Also
\[
 E_+(A)=\frac1p\sum_q|\widehat{1_A}(q)|^4.
\]
Apply \eqref{eq:cayley-fourier-max} to one factor of the nonzero-frequency
fourth moment and Parseval to the other.  This proves
\eqref{eq:cayley-fourier-energy}; the hybrid estimate follows from
Theorem~\ref{thm:czu-cayley}.
\end{proof}

\begin{remark}
This is a genuine dense-regime improvement, not the missing sparse theorem.
The required inversive subgroup sums below the square-root scale are not
controlled by this complete mixed-Weil argument.  In particular the centered
$K^{2+o(1)}$ target remains open in the arithmetic regime relevant to the
present BGS route.  The maintained Python and independent C++ exact
regressions check 807 subgroup designs through $p\le509$, including 364
strict dense-side cases $K^5>p^3$ (certificate
\path{cayley_fourier_hybrid_20260907.json}, indexed in
\path{evidence/provenance.md}).
\end{remark}
